\section{Routing with team constraints}
\label{sec:routing}

A demand specifies two distinct vertices to be joined.  Demands are indexed,
so several demands may specify the same pair.  Besides edge-disjointness of
all routes, we will need internal vertex-disjointness for specified groups
of routes.  The following lemma supplies both requirements.  Its proof uses
the hypergraph formulation and expansion-to-connectivity method of
Buci\'c and Montgomery~\cite[Proposition~8 and Lemma~9]{BM2024}.
The vertex tokens introduced in the proof impose a separate
internal-disjointness constraint for each team; different teams
compete only for graph edges.

\begin{lemma}[Routing]\label{lem:routing}
Let $P$ be a graph of order $r\ge 2$ with cut expansion $h>0$, and put
\begin{equation}\label{eq:routing-radii}
  g=\ceil{8(r/h)\log(2r)},\qquad
  \ell=2g\bigl(1+\ceil{\log r}\bigr).
\end{equation}
Let $\mathcal D$ be a finite set of demands, where demand $p$ has distinct
endvertices $x_p,y_p\in V(P)$.  Suppose that each vertex occurs as an
endvertex in at most $k$ demands, where $k\ge 1$.  Partition $\mathcal D$
into teams of size at most $b$, where $b\ge 1$, and give each demand a set
$Z_p\subseteq V(P)$ of at most $z$ vertices, where $z\ge 0$.
If
\begin{equation}\label{eq:routing-condition}
  h\ge 64\bigl(\ell^2(k+b)+z\bigr),
\end{equation}
then there are simple $x_p$--$y_p$ paths $Q_p$ in $P$, one for each
$p\in\mathcal D$, with the following properties:
\begin{enumerate}[label=\textup{(\roman*)}]
  \item each $Q_p$ has length at most $\ell$ and no internal vertex in $Z_p$;
  \item the paths $Q_p$ are pairwise edge-disjoint;
  \item the internal vertex sets of paths in the same team are pairwise
    disjoint.
\end{enumerate}
The sets $Z_p$ may contain $x_p$ or $y_p$: their restrictions concern only
internal vertices.  If a route must also avoid another demand's endvertices
internally, those vertices may be included in its set $Z_p$.
\end{lemma}

\begin{proof}
There is nothing to prove when $\mathcal D$ is empty.  Suppose it is
nonempty, and write $c(p)$ for the team of demand $p$.

\paragraph{A hypergraph obstruction.}
Take a resource set that is the tagged disjoint union of $E(P)$ and the
tokens $(v,c)$, where $v\in V(P)$ and $c$ is a team.  For each demand $p$,
form a hypergraph $\mathcal H_p$ on this resource set.  Its hyperedges
correspond to all simple $x_p$--$y_p$ paths of length at most $\ell$ whose
internal vertices avoid $Z_p$.  A candidate path contributes its graph
edges and the token $(v,c(p))$ for each of its internal vertices.  A path
of length $t\ge 1$ therefore contributes $2t-1\le 2\ell$ resources.
In particular all hyperedges are nonempty, and all the hypergraphs are
finite.  A choice of pairwise disjoint representatives from the
$\mathcal H_p$ gives exactly the three required properties.

We verify the sufficient condition of \Cref{thm:aharoni-haxell} with
$s=2\ell$.  Suppose, for a contradiction, that it fails for a nonempty
$I\subseteq\mathcal D$.  Let $\mathcal M$ be a maximum matching in
$\bigcup_{p\in I}\mathcal H_p$.  Then
$\abs{\mathcal M}\le 2\ell(\abs I-1)$.  Let $F\subseteq E(P)$ be the
graph-edge resources used by $\mathcal M$, and put
\[
  Y_c=\{v\in V(P):(v,c)\text{ is used by }\mathcal M\}.
\]
Each matching member uses at most $\ell$ graph edges and at most $\ell$
tokens, so
\begin{equation}\label{eq:routing-obstruction-budgets}
  \abs F\le 2\ell^2\abs I,\qquad
  \sum_c\abs{Y_c}\le 2\ell^2\abs I.
\end{equation}
These estimates allow several matching members to come from one demand;
no restriction on their demand indices is being imposed.

For each team let $m_c=\abs{\{p\in I:c(p)=c\}}\le b$.  Thus
\[
  \sum_{p\in I}\abs{Y_{c(p)}}
  =\sum_c m_c\abs{Y_c}
  \le 2b\ell^2\abs I.
\]
Consequently the set
$I_1=\{p\in I:\abs{Y_{c(p)}}\le 4b\ell^2\}$ has size at least
$\abs I/2$.  Choose an inclusion-maximal set $I_2\subseteq I_1$ whose
endpoint pairs are pairwise vertex-disjoint.  Every demand in $I_1$ meets
one of these selected pairs.  The two vertices of any selected pair
occur in at most $2k$ demands in total, so
\begin{equation}\label{eq:routing-pair-count}
  \abs{I_2}\ge\frac{\abs{I_1}}{2k}
  \ge\frac{\abs I}{4k}.
\end{equation}
We will also use the global endpoint count
\begin{equation}\label{eq:routing-endpoint-count}
  2\abs I\le kr.
\end{equation}

\paragraph{A small ball for each selected demand.}
For $p\in I_2$, define
\[
  W_p=(Z_p\cup Y_{c(p)})\setminus\{x_p,y_p\},\qquad
  P_p=(P-F)-W_p.
\]
Here the first deletion removes edges and the second removes vertices.
Both endpoints remain in $P_p$.  By \eqref{eq:routing-condition},
\begin{equation}\label{eq:routing-forbidden-budget}
  \abs{W_p}\le z+4b\ell^2
  \le 4\bigl(z+(k+b)\ell^2\bigr)
  \le h/16\le h/4.
\end{equation}
Set $T=\ell/2=g(1+\ceil{\log r})$, an integer.  If the balls of
radius $T$ about both $x_p$ and $y_p$ in $P_p$ had size greater than
$r/2$, they would intersect.  Joining their paths to an intersection
vertex gives an $x_p$--$y_p$ walk of length at most $2T=\ell$; deleting
closed subwalks gives a simple path of no greater length.  It uses no
edge of $F$ and no vertex of $W_p$.  Its internal vertices are also
different from $x_p,y_p$, and hence avoid both $Z_p$ and $Y_{c(p)}$.
Its hyperedge in $\mathcal H_p$ is therefore disjoint from $\mathcal M$,
contrary to maximality of that matching.

For every $p\in I_2$ we may consequently choose an endpoint $s_p$ whose
radius-$T$ ball in $P_p$ has size at most $r/2$.  Let
$S=\{s_p:p\in I_2\}$.  These roots are distinct because the selected
endpoint pairs are disjoint.  Each root retains its demand label: write
$p(s)$ for the unique $p\in I_2$ with $s=s_p$.  In particular,
\begin{equation}\label{eq:routing-small-root-balls}
  \abs{\{v:\operatorname{dist}_{P_{p(s)}}(s,v)\le T\}}\le r/2
  \qquad(s\in S).
\end{equation}

\paragraph{Growth with different forbidden sets.}
The deleted-edge budget may be too large to grow a ball from one root.
We first grow unions of balls, retaining each root's graph, and then
alternate halving the root set with further growth until one root remains.
For a nonempty $A\subseteq S$ and an integer $t\ge 0$, put
\[
  B_t(A)=\bigcup_{s\in A}
  \{v\in V(P_{p(s)}):\operatorname{dist}_{P_{p(s)}}(s,v)\le t\}.
\]
These are unions of balls in their respective graphs; the forbidden set
is kept with its root throughout.  We claim that
\begin{equation}\label{eq:routing-union-growth}
  \abs{B_{t+1}(A)}
  \ge\left(1+\frac{h}{2r}\right)\abs{B_t(A)}
\end{equation}
whenever
\begin{equation}\label{eq:routing-growth-threshold}
  0<\abs{B_t(A)}\le r/2,
  \qquad \abs F\le h\abs{B_t(A)}/4.
\end{equation}

To prove this, let $B=B_t(A)$ and assign to each $v\in B$ one root
$s(v)\in A$ that reaches $v$ within $t$ steps in $P_{p(s(v))}$.
Consider a boundary edge $vw$ of $B$, with $v\in B$ and $w\notin B$.
It extends this particular root's path whenever $vw\notin F$ and
$w\notin W_{p(s(v))}$.  In that event $w\in B_{t+1}(A)\setminus B$.
Because $P$ is simple, at most $\abs{W_{p(s(v))}}\le h/4$ boundary
edges at $v$ fail the latter condition.  Cut expansion and
\eqref{eq:routing-growth-threshold} therefore leave at least
\[
  h\abs B-(h/4)\abs B-\abs F\ge(h/2)\abs B
\]
usable boundary edges.  Each new vertex receives at most $r$ such
edges.  Thus at least $h\abs B/(2r)$ new vertices are reached, proving
\eqref{eq:routing-union-growth}.

It follows that, for every integer $t\ge 0$,
\begin{equation}\label{eq:routing-growth-block}
  \abs{B_t(A)}\ge 1\ \text{ and }\quad
  \abs F\le h\abs{B_t(A)}/4
  \quad\Longrightarrow\quad
  \abs{B_{t+g}(A)}>r/2.
\end{equation}
Indeed, the conclusion is immediate by monotonicity if the starting
union already has size greater than $r/2$.  Otherwise, if it still had
size at most $r/2$ after $g$ more steps, all intervening unions would
have size at most $r/2$, while their monotonicity would preserve the
edge-deletion threshold.  We could apply
\eqref{eq:routing-union-growth} at all $g$ steps.  Singleton cuts in
the simple graph $P$ give $0<h\le r-1<r$, and the elementary inequality
$\log(1+h/(2r))\ge h/(4r)$ then yields
\[
  \abs{B_{t+g}(A)}
  \ge \left(1+\frac{h}{2r}\right)^g\abs{B_t(A)}
  \ge 2^{gh/(4r)}
  \ge (2r)^2>r/2,
\]
a contradiction.  This also verifies the growth time with the ceiling
in the definition of $g$.

\paragraph{Restarting after halving the roots.}
Initially $B_0(S)=S$.  By \eqref{eq:routing-obstruction-budgets},
\eqref{eq:routing-pair-count}, and \eqref{eq:routing-condition},
\[
  \abs F\le 2\ell^2\abs I
  \le 8k\ell^2\abs S\le h\abs S/4.
\]
The growth block \eqref{eq:routing-growth-block} gives
$\abs{B_g(S)}>r/2$.

Suppose now that a nonempty root set $A$ has
$\abs{B_t(A)}>r/2$ and $\abs A\ge 2$.  Partition $A$ into two
nonempty sets whose sizes differ by at most one.  Their radius-$t$
ball unions together cover $B_t(A)$, so one of them, say $A'$, satisfies
\[
  \abs{A'}\le\ceil{\abs A/2},\qquad
  \abs{B_t(A')}>r/4.
\]
The original matching has a bound independent of the number of retained
roots: by \eqref{eq:routing-obstruction-budgets},
\eqref{eq:routing-endpoint-count}, and \eqref{eq:routing-condition},
\begin{equation}\label{eq:routing-restart-budget}
  \abs F\le 2\ell^2\abs I
  \le \ell^2kr\le hr/64
  < h\abs{B_t(A')}/4.
\end{equation}
Thus another growth block gives $\abs{B_{t+g}(A')}>r/2$.

Start with $A=S$ at radius $g$ and repeat this operation until only one
root remains.  After $j$ halvings the number of roots is at most
$\ceil{\abs S/2^j}$; this follows inductively from rounding each half
upward.  Hence at most $\ceil{\log\abs S}$ halvings are needed,
including zero if $\abs S=1$.  At termination a single root has a ball
larger than $r/2$ at an integer radius at most
\[
  g\bigl(1+\ceil{\log\abs S}\bigr)
  \le g\bigl(1+\ceil{\log r}\bigr)=T.
\]
This contradicts \eqref{eq:routing-small-root-balls} and monotonicity of
that root's balls.

There is therefore no set $I$ for which the matching condition fails.
Applying \Cref{thm:aharoni-haxell} gives disjoint hypergraph
representatives, and their paths prove the lemma.
\end{proof}

For later estimates, the exact radii also give the useful uniform bound
\begin{equation}\label{eq:routing-length-bound}
  \ell\le 108(r/h)(\log r)^2
  =O\bigl((r/h)(\log r)^2\bigr).
\end{equation}
Indeed, $h<r$ and $r\ge 2$ imply
$g\le 9(r/h)\log(2r)$,
$\log(2r)\le 2\log r$, and
$1+\ceil{\log r}\le 3\log r$.
